Everything is computed in Gauge Theory Gravity (Definition \claimref{obj.synthesis.gtg-flat-spacetime-gauge-fields}) and in the conformal model $\Cl(2,4)$. At $p$ the fibre is the background vectors; the physical vector over $x$ is $\hb^{-1}(x)$, the physical covector over a background gradient $\xi$ is $\hb^{T}(\xi)$, and $Q_p(x) = Q(\hb^{-1}(x))$, $Q^{*}_p(\xi) = Q(\hb^{T}(\xi))$. In the static gauge $\hb^{-1}(e_{t}) = \sqrt f\,\gamma_{0}$, $\hb^{-1}(e_{r}) = \gamma_{1}/\sqrt f$, $\hb^{-1}(e_{\theta}) = r\gamma_{2}$, $\hb^{-1}(e_{\phi}) = r\sin\theta\,\gamma_{3}$ on a constant orthonormal physical frame, so $Q_p$ has the Gram $\operatorname{diag}(f, -f^{-1}, -r^{2}, -r^{2}\sin^{2}\theta)$ displayed in the statement. Write a vector of $\Cl(2,4)$ as $N = s\,n_{0} + X + t\,\nInf$ with $X$ a physical vector and the null pair normalised by $n_{0}\cdot\nInf = -1$ (here $n_{0} = (e_{-} - e_{+})/2$, $\nInf = e_{-} + e_{+}$); these are the triples $(s, X, t)$ of \claimref{transport.swap-rescale-normal-form}, $N^{2} = Q(X) - 2st$, and the lift of $x$ is $X_p(x) = n_{0} + \hb^{-1}(x) + \tfrac12 Q_p(x)\,\nInf$, the triple $(1, \hb^{-1}(x), \tfrac12 Q_p(x))$.

(1) $X_p(x)^{2} = Q_p(x) - 2\cdot\tfrac12 Q_p(x) = 0$: the incidence condition $Q_p(x) = 2st$ holds for every $x$. The transported reciprocal is the reflection in $e_{+}$, $\iota[N] = [e_{+}Ne_{+}]$: since $e_{+}\cdot n_{0} = -\tfrac12$, $e_{+}\cdot\nInf = 1$ and $e_{+}$ is orthogonal to spacetime, $e_{+}Ne_{+} = -(2t\,n_{0} + X + \tfrac{s}{2}\,\nInf)$, which is the triple $(2t, X, \tfrac{s}{2})$ up to the scale $-1$. As a versor it preserves the quadric, $(e_{+}Ne_{+})^{2} = e_{+}N^{2}e_{+} = N^{2}$ ($N^{2}$ is a scalar and $e_{+}^{2} = 1$), so it preserves incidence, and $e_{+}(e_{+}Ne_{+})e_{+} = N$: it is an involution.

(2) $\iota(X_p(x)) = [\,Q_p(x)\,n_{0} + \hb^{-1}(x) + \tfrac12\nInf\,]$. If $Q_p(x) \neq 0$, the read-out $\dsc(N) = X/s$ gives $\hb^{-1}(x)/Q_p(x)$, the physical vector over $x/Q_p(x)$, the classical reciprocal $\rho(x)$ of \cite[Lem.~4.23(2)]{ButlerT1} (\claimref{transport.descent-read-out-map}).

(3) For $x = e_{t} + f\,e_{r}$, $\hb^{-1}(x) = \sqrt f\,(\gamma_{0} + \gamma_{1})$, a physically null vector, so $Q_p(x) = 0$, $X_p(x) = n_{0} + \hb^{-1}(x)$ and $\iota(X_p(x)) = [\,\hb^{-1}(x) + \tfrac12\nInf\,]$: its $n_{0}$ slot vanishes, which is the defining condition of an ideal point, and $\dsc$ requires division by it. This is the ideal value of \cite[Thm.~4.10(3)]{ButlerT1} (\claimref{transport.totality}).

(4) The one-forms of the statement are background gradients: $dr = \nabla r = e^{r} = -e_{r}$ and $dt_{\mathrm{PG}} = \nabla t_{\mathrm{PG}} = e^{t} = e_{t}$ in the Painlev\'e-Gullstrand labels, so the generator of the radical line of \cite[Thm.~8.7]{ButlerT4} (\claimref{sheaf.schwarzschild-degeneration-boundary-tame}) is the background vector $\xi = \sqrt{2M}\,e_{t} - \sqrt r\,e_{r}$. In the stored gauge $\hb = 1 - \beta P$ the adjoint is $\hb^{T}(a) = a - \beta\,(a\cdot e_{r})\,e_{t}$ (Section 11 of the dictionary, $\beta = \sqrt{2M/r}$), and $\xi\cdot e_{r} = \sqrt r$, so
\[
\hb^{T}(\xi) = (\sqrt{2M} - \beta\sqrt r)\,e_{t} - \sqrt r\,e_{r} = -\sqrt r\,e_{r}, \qquad Q^{*}_p(\xi) = Q(-\sqrt r\,e_{r}) = -r,
\]
using $\beta\sqrt r = \sqrt{2M}$: exactly, at every radius, the horizon included, with no pole. The physical covector $\hb^{T}(\xi) = -\sqrt r\,e_{r}$ is radial outright, so the corank-one radical of the completion and the dying direction $e_{r}$ of Step 1, where Propositions \claimref{singularity.cones-fold} and \claimref{singularity.degeneration-exactly} fold the cones in the same chart, are the same direction at every radius, not merely in the limit. In the static labels the same covector reads $\xi = \sqrt{2M}\,\nabla t_{s} + (\sqrt r/f)\,\nabla r$, because $t_{\mathrm{PG}} = t_{s} + \int\beta f^{-1}dr$ gives $\nabla t_{\mathrm{PG}} = \nabla t_{s} + (\beta/f)\nabla r$ and $\sqrt r\,(1 + \beta^{2}/f) = \sqrt r/f$ by $f + \beta^{2} = 1$; with the static gauge, $\hb^{T}(\nabla t_{s}) = \gamma^{0}/\sqrt f$ and $\hb^{T}(\nabla r) = \sqrt f\,\gamma^{1}$, so $Q^{*}_p(\xi) = 2M/f - f\cdot r/f^{2} = (2M - r)/f = -r$ again. (The Killing-time gradient alone has $Q^{*}_p(\nabla t_{s}) = 1/f$, with a pole at the horizon where nothing degenerates; it spans the radical line only on $E$, \cite[Thm.~8.7(4)]{ButlerT4}; $\xi$ is the regular generator.) Hence $\iota(X_p(\xi)) = [\,-r\,n_{0} + \hb^{T}(\xi) + \tfrac12\nInf\,]$ has nonvanishing $n_{0}$ slot at every $r > 0$, where $\dsc$ returns $-\hb^{T}(\xi)/r$, the physical covector over $-\xi/r$; as a polynomial in $u = \sqrt r$ its slots are $(-u^{2}, -u\,e_{r}, \tfrac12)$, whose value at $u = 0$ is the ideal point $[\tfrac12\nInf]$. That the failure locus is exactly $E$, and that no shadow wall forms elsewhere, is \cite[Thm.~8.12]{ButlerT4} (\claimref{sheaf.one-sided-completion-singularity}).

The check verifies (1) to (3) blade by blade in $\Cl(2,4)$ at three static rational points (lift nullity, $e_{+}Ne_{+}$ and its slots, the involution, the quadric value of a non-incident $N$, the reciprocal and the ideal point with the division failing), and (4) with the shared engine's Painlev\'e-Gullstrand and static gauge fields: $\hb^{T}(\xi) = -\sqrt r\,e_{r}$ and $Q^{*}_p(\xi) = -r$ at three masses and four radii inside, on and outside the horizon, $-r$ again in the static gauge with $Q^{*}_p(\nabla t_{s}) = 1/f$, and the read-out $-\hb^{T}(\xi)/r$.

\emph{In components} (a labelled reading). In the static cometric $(g_p^{-1})^{tt} = 1/f$, $(g_p^{-1})^{rr} = -f$ the computation of (4) is $Q^{*}_p(\xi) = 2M/f - f\,(r/f^{2}) = -r$; in the Painlev\'e-Gullstrand cometric $g^{tt} = 1$, $g^{tr} = -\beta$, $g^{rr} = \beta^{2} - 1$ the components of $k^{\sharp}$ are $\sqrt{2M} - \beta\sqrt r = 0$ and $-\beta\sqrt{2M} + (\beta^{2} - 1)\sqrt r = -\sqrt r$, which is $\hb^{T}(\xi) = -\sqrt r\,e_{r}$; and the polynomial Gram matrix of \cite[Thm.~8.7]{ButlerT4} gives $g_{k} = (m^{2} - r) - 2m^{2} + m^{2} = -r$ with $m = \sqrt{2M}$, that theorem's clause (3). The identities of (1) to (4) are also machine-verified in the series' legacy script (parts L1 to L14).
